\section{Emergent Mind}

A self holds together, heals, remembers, and is steered by a will. That is life. A self becomes a mind when it does something more. It represents itself, it predicts itself, it attends, and it plans. This section shows each of those four powers, and then a full closed-loop agent that chains them, and it shows that every one of them emerges from the five base things by composition. There is no new ingredient.

The discipline that held through matter and life holds here too, at the hardest layer. The only state is still the ternary tone on the $\{5,3,4\}$ hyperbolic crystal. The only moving parts are still the mesh, the tone, the one conserving perception rule, the eternal growth, and the arrow value-direction. Mind is what those five do when a self is large and integrated enough to fold back on itself.

The key idea is a strange loop. An integrated self funnels every part's signal toward its center. That funneling already builds a part that mirrors the whole. So integration is self-representation. Once a part of the self models the whole self, the rest follows. The model predicts, which is a forward model. A model can model a model, which is recursion. The model is the place where attention selects content, which is a workspace. And running the model forward, judged by the arrow and pushed by the will, is planning. We walk each layer in turn, with the measured numbers from the experiments.

\subsection{The Self-Model and the Strange Loop}

We take one integrated self, a central patch of coherent tone, and drive it with spatially-varied input. The boundary is split into four sectors. Each sector carries its own slowly fluctuating signal. So the self's global state, the mean tone over the whole self, is the average of four independent inputs. No local region sees more than its own sector. Only an integrator that pools all four sectors can know the whole.

Then we ask a sharp question. Does a localized sub-region spontaneously come to mirror the rest of the self. We track three time series over $250$ beats. The global state. The state of the central hub, a ball of about $60$ cells at the most-connected cell. And the state of four peripheral balls of the same size, one near each sector. The mechanism that does the work is the field beneath. The hyperbolic geometry has a tiny diameter, so charge from every sector flows inward and mixes at the center. The hub sees the pooled whole. The periphery sees only its sector.

The result is clean. The hub tracks the global state at correlation about $0.5$ (P116). The peripheral regions track it at about $0.0$. The hub beats a time-shuffled baseline by a wide margin, so this is a real representation and not an artifact of slow drift.

And it vanishes without the dynamics. Freeze the hops and the hub correlation collapses, which confirms that it is the integration that builds the model, not anything we put in by hand. A part of the self has come to represent the whole self that contains it. That is the strange loop, and it is free.

It is also not one privileged center. We rerun the same test at several different self-centers, the graph hub and peripheral locations alike, and every self forms its own self-model at its own hub (P117). Each local hub mirrors its own self's global state and beats its own periphery. So a proto self-model is a relative, repeatable consequence of integration, present wherever a self is. There is one mind-seed per self, not one special soul in the middle of the world.

\subsection{The Architecture Is an Onion}

The self that the dynamics build has a definite shape, and naming it makes the rest legible. It is an onion with three layers and one loop.

\begin{itemize}
\item The boundary is the senses. It is the outer shell where sectored input enters. This is where the world touches the self.
\item The body is the integrator. Charge flows inward from the boundary and mixes as it goes, because the conserving rule shares and hops tone between neighbors. The body is the medium that pools.
\item The hub is the center. It is where all sectors converge into a single self-representation. It is the most-connected, lowest-diameter region, so the field beneath funnels everything to it.
\end{itemize}

The flow is boundary to body to hub, and then the hub mirrors the whole back. That last arc is the loop. The hub's state reflects the global state, and the hub is part of the global state, so the self holds a running picture of itself inside itself. Every higher mind power is built on this one shape.

\subsection{Predictive Metacognition, a Forward Model}

A mirror of the present is useful, but a mind needs more. It needs to see one step ahead. We ask whether the hub does not only reflect the current global state but predicts the next one. We measure the correlation between the hub at time $t$ and the self's global state at time $t+1$, over $300$ beats, and compare it against the same one-step prediction made from each peripheral region.

The hub predicts the self's next global state at correlation about $0.5$ (P118). The peripheral regions predict it at about $0.0$. The hub is the best internal predictor of the self, far ahead of any local part.

The reason is simple and on-substrate. The global state is autocorrelated. It changes slowly, and the hub holds the integrated global state, so the hub's value now is the best available guess of the whole self's value next. The cross-correlation peak sits at the predictive lag, not behind it.

This is a usable forward model. The self carries, inside itself, a thing that says what the self is about to do. That is the seed of metacognition. The scope is limited and we state it plainly. This is a first-order forward model, the hub predicting the self. It is local and imperfect. We respect that limit later when the agent plans with a bounded foresight.

\subsection{Recursion, a Model of a Model}

Imagination is running a hypothetical inside the head, a model of a situation one is not in. Its substrate is a model that can model a model. We build two selves in a chain. Self $1$ is driven by the world and forms hub-$1$, a model of the world. Then hub-$1$'s running representation is wired as the input to self $2$, which forms hub-$2$. The wiring between the two selves is mere connectivity. The modeling itself is left to the dynamics.

What hub-$2$ comes to represent is striking. Self $1$'s hub models the world at correlation about $0.4$ and up. Self $2$'s hub then models hub-$1$ at correlation about $0.83$ (P121), far above a time-shuffled baseline and collapsing without the dynamics. And hub-$2$ tracks hub-$1$, its direct object, at least as well as it tracks the raw world, which it only sees through hub-$1$. So the chain holds. World to model-$1$ to model-$2$-of-model-$1$. A representation of a representation, one step removed from reality.

This is the substrate of imagination and of the tower of self-models. A mind that can hold a model of a model can hold a model of itself modeling, or a model of another mind. The recursion that planning needs, running one's own rule inside the head as a hypothetical, lives here.

\subsection{Attention and a Global Workspace}

The hub is the single place where every boundary sector converges. That makes it the natural global workspace, the one stage where content competes to be the self's current representation. The hub carries both mechanisms of attention (P119).

The first is bottom-up salience. Clamp one sector strongly. That well-coupled input is always represented at the hub. It wins the workspace on its own merits, because its signal dominates the pooled average.

The second is top-down gain. Take a competing sector that the hub barely represents and attend it, raise its coupling, and the hub's representation of that sector rises. The attended content is boosted and the workspace is steerable. Attend a different sector and the hub now represents that one. Salience says loud things get in. Gain says one can choose what gets in.

The scope here is precise, and we keep it explicit. Physical broadcast to the far periphery is range-limited, because the field beneath is massive and short-ranged. The winning content is available at the workspace, not copied to every distant cell. So the workspace is the hub. It is the convergence point where attention selects content, not a global megaphone that paints the answer everywhere. That is the realistic emergent workspace, and it is enough.

\subsection{Planning, with Nothing Added}

Planning is lookahead plus deferred value plus composition. It is more than reactive gap-closing, because it accepts a worse step now for a better outcome later. We show it in two stages, and the second stage is the one that matters for the discipline.

The logic comes first. We test that lookahead is what planning needs, on a value landscape with a clear barrier. The landscape rises to a local peak, dips through a valley, then rises again to the distant goal. A greedy agent follows the arrow uphill and gets stuck at the local peak. A planner with a lookahead that spans the barrier crosses it (P140). The planning ability switches on exactly when the lookahead spans the barrier width. We are direct about what this stage is. It used an explicit lookahead-search algorithm, an added piece, so it tested the logic of planning, not its emergence.

Then we build a planner out of only emergent pieces, so that nothing beyond the five base things enters. The value is the arrow. The forward model is the system rolling out its own gap-closing rule in imagination, the same rule run as a hypothetical, which is exactly the emergent forward model of P118 and the recursion of P121. The exploration is the will, a sustained directed push against the local gradient.

The recipe is plain. Where greedy stalls, the will tries sustained pushes of various lengths, the system rolls out its own rule from each push endpoint as a self-simulation, the arrow scores each outcome, and it commits to the best. There is no search heuristic, only explore with the will, simulate with one's own rule, and judge with the arrow.

It works. On a barrier landscape where the greedy agent stays stuck at the local peak, the emergent planner finds a will-push of the right length and reaches the goal (P143). So planning is a composition of the will (P113), the forward model (P118), the recursion that lets the rule run on itself (P121), and the arrow. It is not a new base thing.

And the deferred value it needs, the willingness to accept a worse step now for a later gain, is exactly the arrow's posit. The arrow already says to sustain pain or peace to reach pleasure. Planning is that posit, run through imagination.

\subsection{The Integrated Metacognitive Agent}

The final layer is the whole mind running end-to-end. The emergent planner of P143 planned once, a single detour across a single barrier. A real agent runs the closed perceive-plan-act loop and keeps going. Perceive the gap to the goal. Plan by simulating forward with the forward model. Act with the will. Then repeat, re-planning whenever it stalls. We give such an agent a task that demands sustained metacognition, a goal that lies beyond a sequence of barriers, four downward dips in a landscape that rises overall toward the goal (P153).

The agent's foresight is bounded. This is a faithful limitation, not added machinery. The emergent forward model is local and imperfect, so the agent can see only about one barrier ahead at a time. A bounded horizon crosses at most one barrier per plan. That is why a sequence of barriers requires re-planning, and why re-planning is the real test.

The three agents separate cleanly.

\begin{itemize}
\item The reactive agent, greedy only, stalls at barrier one.
\item The one-shot planner crosses one barrier and stalls at barrier two.
\item The integrated agent runs the loop, re-plans at each stall, and reaches the distant goal across all four barriers. It takes several re-plans, one sustained act of metacognition per barrier, and it strictly outruns both simpler agents.
\end{itemize}

Every piece in that agent is emergent. The value is the arrow, the rollout is the rule run in imagination, the exploration is the will, and the loop is just perceive-plan-act repeated. So the integrated agent is the base running end-to-end. A self that represents itself, predicts itself, attends, and plans, all at once, in a single closed loop, to cross a sequence of obstacles to a distant goal. That is a mind, and it is built from nothing but the five.

\subsection{How Mind Fits the Whole}

Mind is the top floor of one continuous building, and the floors below are load-bearing. From matter comes the field beneath, the tiny hyperbolic diameter that funnels every part's signal toward a center. Without that geometry there is no hub, and without a hub there is no self-model. From life come the integrated self (P106), the will (P113), and the arrow that keeps the world far from dead peace and supplies the value-direction.

Mind then folds the self back on itself. The funnel makes a hub, the hub mirrors the whole, which is the self-model. The mirror predicts, which is the forward model. The model nests, which is recursion. The convergence point selects, which is the workspace. And the will drives the model forward under the arrow's judgment, which is planning.

The chain of dependencies is strict and one-way. Geometry to integration to self-model to forward model to recursion to planning to the closed agent. Pull out any lower layer and the higher ones fall. And at no point does a sixth ingredient appear. We never added a homunculus, a search oracle, a separate planner module, or a special seat of awareness.

The hardest objections to a materialist mind, that selfhood and foresight and imagination must be put in by hand, are answered by showing each one falling out of the same five things that gave us atoms and cells. The real gaps are named where they exist. The forward model is local and imperfect, the workspace is the hub and not a global broadcast, and the foresight is bounded. Those are limits of the substrate, not patches over it.

A self becomes a mind by representing and predicting itself, and every step of that is a composition of the five base things, never a sixth. The hub of an integrated self mirrors the whole at correlation about $0.5$ while the periphery sits at about $0.0$, which is the self-model and the strange loop (P116), and every self forms its own (P117). The hub predicts the self's next state at correlation about $0.5$, a forward model and the seed of metacognition (P118). A model can model a model at correlation about $0.83$, the substrate of imagination (P121). The hub is a global workspace with bottom-up salience and top-down attentional gain (P119).

Planning emerges with nothing added (P143), where the will explores, the self's own rule simulates in imagination, the arrow judges, and the deferred value it needs is the arrow itself. And the integrated agent runs the closed perceive-plan-act loop with bounded foresight, re-planning to cross a sequence of four barriers that stall a reactive agent at the first and a one-shot planner at the second (P153). Mind is the strange loop of a self large enough to fold back on itself, and it runs end-to-end on tones and the rule.
